\section{Reserve bounds and threshold reductions}\label{sec:versions}

Two further parts of the first-passage argument depend on retaining the
appropriate finite-scale information: lifting a congruence between offset
integers requires a size bound, and passing to logarithmic density requires
the finite dyadic weights and the correct range of a threshold function.

\subsection{A reserve for \texorpdfstring{$3$}{3}-adic separation}

The $3$-adic separation argument in \cite[Section~6]{Tao7} retains a
reserve of $0.99C^2\log n$. The following estimate gives an explicit
sufficient margin for this size argument. Put
\[
 a=\log(4/3),\qquad b=\log2,\qquad \rho_*=\frac{b}{4a}.
\]
For positive integers $a_1,\ldots,a_k$, write $S_j=a_1+\cdots+a_j$
and $l=S_k$.  Clearing the dyadic denominator of the reversed offset
produces the positive integer
\[
 U=\sum_{j=1}^k3^{j-1}2^{l-S_j}.
\]
Bounding $U$ below $3^n$ is what permits an equality modulo $3^n$
between two such integers to be lifted to an equality of integers.

\begin{proposition}[An explicit reserve bound]\label{prop:reserve}
Suppose $n\geq2$, $C>0$, and $\rho>\rho_*$ satisfy
\[
 S_j\geq2j-2-C(\sqrt{j\log n}+\log n)\quad(1\leq j\leq k),
 \qquad
 l\leq n\frac{\log3}{b}-\rho C^2\log n.
\]
Put $\gamma=\rho b-b^2/(4a)>0$.  Then
\begin{equation}\label{eq:reserve}
 \frac{U}{3^n}\leq\frac83 n^{1-\gamma C^2+bC}.
\end{equation}
In particular $U<3^n$ if $\gamma C^2-bC\geq3$.
\end{proposition}
\begin{proof}
For each $j$, the hypotheses give
\[
 \frac{3^{j-1}2^{l-S_j}}{3^n}
 \leq\frac43\exp\left(-\rho b C^2\log n+bC\log n
                  -aj+bC\sqrt{j\log n}\right).
\]
Complete the square:
\[
 -aj+bC\sqrt{j\log n}
 =-a\left(\sqrt j-\frac{bC\sqrt{\log n}}{2a}\right)^2
  +\frac{b^2C^2\log n}{4a}.
\]
Every summand is thus at most $\frac43 n^{-\gamma C^2+bC}$.
Since $k\leq l<n\log3/\log2<2n$, summation gives
\eqref{eq:reserve}.  Under the stated last condition its right-hand
side is at most $8/(3n^2)\leq2/3<1$.
\end{proof}

Numerically $\rho_*=0.602355\ldots$, so the coefficient $0.99$ used by
Tao satisfies the strict inequality required here. The displayed bound
quantifies the sufficient reserve; it does not assert that this threshold
is necessary for other estimates of the same sum.

Here is the connection with the stopping event used there.  If
$T=n\log3/\log2-C^2\log n$ and the accumulated valuation first exceeds
$T$ at index $k+1$, the source concentration condition for the last two
coordinates gives
\[
 S_{k+1}\leq T+1+C(\sqrt{\log n}+\log n)
 \leq T+6C\log n.
\]
For $C\geq600$ this is at most
$n\log3/\log2-0.99C^2\log n$, exactly the retained reserve.
In deriving the first inequality, use $S_k\leq T$ and
$a_k\geq1$ in
$|a_k+a_{k+1}-2|\leq C(\sqrt{\log n}+\log n)$.
The source has already restricted the index to
$k=n\log3/(2\log2)+O(C\sqrt{n\log n})$, so $k\geq1$ for sufficiently
large $n$ at fixed $C$.  The prefix bound in
Proposition~\ref{prop:reserve} allows the additive two arising from
the source's interval indexing: for $j\geq2$, its interval starting
at index one gives $S_j\geq2(j-1)-C(\sqrt{(j-1)\log n}+\log n)$;
the case $j=1$ follows from $a_1\geq1$.

The role of the reserve is to ensure that congruent offset integers in
this size range are equal as integers.  Fourier decay additionally uses
the analytic estimates in Tao's argument.

\subsection{Dyadic weights and diverging thresholds}

The density factor in v7 Section~1.2 assigns mass $2^{-a}$ to
$\vnu(n)=a$.  Proposition~\ref{prop:dyadic} shows that the absolute
logarithmic density is $2^{-a-1}$; already $a=0$ distinguishes the
formulas.  The source prints the partial sum as $1-2^{-A_0}$;
the corrected finite partial sum is $1-2^{-A_0-1}$.
These two expressions tend to one, so this slip does not alter the limiting
conclusion, but it must be corrected in a finite-weight calculation.

At the end of Section~3, the source sets
$\widetilde f(X)=\inf_{m\geq X,\ m\text{ odd}}f(m)$ and applies it to a
sum over $m\leq X$.  This lower envelope does not control the earlier
values.  For a counterexample to the displayed bound, take $f(1)=0$ and
$f(m)=e^m$ for odd
$m\geq3$.  The source's bad harmonic sum is exactly $1$, since
$\Syr_{\min}(1)=1$ and $\Syr_{\min}(m)\leq m<e^m$ for all other $m$.
Its proposed bound is $O(\log X/X^c)$ and tends to zero.
The fixed-threshold argument of Corollary~\ref{cor:tao} supplies the
complete conclusion, including the strict inequality and arbitrary
real-valued nonmonotone functions, without that display.

\subsection{Source locations}

The following locators refer to \texttt{collatz.tex} in arXiv v7; the
source manifest records its SHA-256 identity.  Line numbers refer to that
TeX source, alongside the section and proposition numbers, rather than
to the journal pagination.

\begin{center}\small
\begin{tabular}{@{}>{\raggedright\arraybackslash}p{0.25\textwidth}>{\raggedright\arraybackslash}p{0.25\textwidth}>{\raggedright\arraybackslash}p{0.42\textwidth}@{}}
\toprule
Source locus & Estimate & Corresponding result\\
\midrule
Section 1.2, line 206 & Dyadic density factor & Exact finite transport,
Propositions~\ref{prop:dyadic} and \ref{prop:sampling}.\\[4pt]
Section 3, lines 591--593 & Tail envelope used on the wrong range &
Counterexample to the display; full replacement in
Corollary~\ref{cor:tao}.\\[4pt]
Proposition 5.2 proof, lines 759--762 & Additive input buffer &
Multiplicative buffer, Proposition~\ref{prop:buffer}.\\[4pt]
Lemma 5.3 proof, lines 868, 899, 913--914 & Previous-iterate exponent;
nonuniform final sum & Exact affine coefficient and uniformly summable
weights, Proposition~\ref{prop:cn}.\\[4pt]
Section 6, lines 989, 997, 1080--1084 & Retained offset reserve &
Explicit sufficient size bound,
Proposition~\ref{prop:reserve}.\\
\bottomrule
\end{tabular}
\end{center}
